% TEMPLATE-MILC-v3.tex - plantilla maestra UNICA de las guias MILC v3.
%
% La usan las cuatro familias de guias, cada una con su driver:
%   · Grados 6-11          -> scripts/build-guias-g11.py
%   · Bebras               -> scripts/build-guias-bebras.py
%   · Semillero            -> scripts/build-guias-semillero.py
%   · Territorio interior  -> scripts/build-guias-territorio-interior.py
%
% Antes habia tres copias casi identicas (~400 lineas iguales, 6 distintas):
% cada mejora habia que aplicarla tres veces, y en la practica no se hacia
% (el motor de recursos vivia solo en dos de ellas). Lo que varia entre
% familias esta parametrizado en 6 placeholders que cada driver llena
% (se nombran aqui SIN su sintaxis real: el builder reemplaza por texto plano,
% tambien dentro de los comentarios, y un valor multilinea rompe el preambulo):
%
%   HEADER_IZQ        encabezado izquierdo de pagina
%   HEADER_DER        encabezado derecho de pagina
%   PORTADA_ETIQUETA  rotulo sobre el numero grande (GRADO/MOMENTO/MODULO)
%   PORTADA_SERIE     linea de serie de la portada
%   PORTADA_ID        identificador de la guia en la portada
%   PORTADA_CIRCULO   el circulito de grado (°), o vacio si no aplica
%   PDF_TITULO        titulo en texto plano (sin \textbf ni comillas TeX) para
%                     los metadatos del PDF (/Title); lo produce lib_xelatex.texto_pdf
%
% Composicion (auditoria 2026-09-03): las bandas de estacion piden espacio
% (\Needspace*) y prohiben el salto justo despues (\nopagebreak), asi no quedan
% huerfanas al pie; las cajas largas (softbox, drawbox, infoband, puentebox) son
% `breakable`, asi una caja de media pagina ya no arrastra todo a la siguiente
% dejando la anterior al 40 %. Todo texto sobre mostaza o verde va en milcNegro
% (blanco sobre #E5B400 da 1,93:1 y sobre #58A923 2,95:1; ninguno pasa WCAG AA).
% El resto de claves las documenta content/guias/_SCHEMA.md. El script valida
% que no queden placeholders sin reemplazar antes de escribir el archivo final.

% Etiquetado PDF (latex-lab): /Lang, estructura y texto alternativo de las
% figuras, para lectores de pantalla. Debe ir ANTES de \documentclass.
\DocumentMetadata{lang=es-CO,tagging=on}
\documentclass[11pt,letterpaper]{article}
% headheight=24pt: la cabecera derecha lleva dos lineas (identidad de la guia
% y estacion actual). headsep se acorta en la misma medida para que el bloque
% de texto quede exactamente donde estaba con headheight=12pt.
\usepackage[margin=2.0cm,headheight=24pt,headsep=13pt]{geometry}
\usepackage[table]{xcolor}
\usepackage{fontspec}
\usepackage[spanish]{babel}
\usepackage{tikz}
% Librerías TikZ para los diagramas de las guías (bloque `recursos:`):
%  · babel  → babel en español vuelve activos caracteres como «>», y TikZ los usa
%             en sus opciones (>=stealth, ->). Sin esta librería, cualquier
%             diagrama con flechas revienta con «\language@active@arg> has an
%             extra }». Debe ir ANTES de que se dibuje nada.
%  · positioning        → sintaxis «below left=2pt and 3pt».
%  · decorations.pathreplacing → llaves ({brace}) para señalar rangos.
\usetikzlibrary{babel,positioning,decorations.pathreplacing}
\usepackage{graphicx}
% Circuitos: el trabajo de electrónica se hace hoy con micro:bit y sus sensores
% integrados (MakeCode, sin cableado), así que no se necesitan esquemáticos.
% Si algún día entran montajes con protoboard, basta descomentar esta línea:
% los diagramas .tex/.tikz declarados en `recursos:` ya se incrustan solos.
% \usepackage{circuitikz}
\usepackage[most]{tcolorbox}
\usepackage{tabularx}
\usepackage{array}
\usepackage{fancyhdr}
\usepackage{titlesec}
\usepackage[document]{ragged2e}  % bandera a la derecha en todo el cuerpo
\usepackage{enumitem}
\usepackage{microtype}
\usepackage{etoolbox}
\usepackage{needspace}
% hyperref al final del preambulo: metadatos del PDF (titulo, autor, idioma,
% familia), marcadores por estacion (\pdfbookmark) y enlaces sin recuadro.
\usepackage[hidelinks,
  pdftitle={Mini estudio de datos — ocho etapas para responder una pregunta del colegio},
  pdfauthor={PhD. Álvaro Cárdenas Orozco},
  pdfsubject={Tecnología e Informática · Grado 8 · Guía 9 · MILC v3},
  pdflang={es-CO},
  bookmarksopen=true,bookmarksnumbered=false]{hyperref}

% Helvetica Neue no trae la flecha U+2192, asi que cada «→» escrito literalmente
% en el YAML se compilaba a NADA, en silencio (462 flechas en 61 guias: rutas del
% tipo «paleta → tipografia → lockup» salian rotas). Mapeamos el caracter a su
% version matematica, que si tiene glifo. Asi el autor puede seguir escribiendo
% «→» comodamente en el YAML.
\usepackage{newunicodechar}
\newunicodechar{→}{\ensuremath{\rightarrow}}
% Mismo problema con los simbolos de marca y vinetas: el «✓» de «una palomita ✓»
% salia como cuadrito vacio en el PDF de todas las guias que lo usaban.
\usepackage{pifont}
\newunicodechar{✓}{\ding{51}}
\newunicodechar{✗}{\ding{55}}
\newunicodechar{✔}{\ding{52}}
\newunicodechar{•}{\textbullet}
\newunicodechar{≥}{\ensuremath{\geq}}
\newunicodechar{≤}{\ensuremath{\leq}}

\setmainfont{Helvetica Neue}
\setsansfont{Helvetica Neue}
\setlength{\parindent}{0pt}
\setlength{\parskip}{6pt}
% Sin particion de palabras: en 6.º-7.º una palabra cortada por guion al final
% de linea («Comuni-cacion») frena la lectura mucho mas que una linea corta.
\hyphenpenalty=10000
\exhyphenpenalty=10000
\pretolerance=2000
\tolerance=3000
\setlength{\emergencystretch}{3em}
\linespread{1.10}
\renewcommand{\arraystretch}{1.25}
\setlist{leftmargin=5mm,itemsep=1mm,topsep=1mm}
\newcolumntype{Y}{>{\RaggedRight\arraybackslash}X}

\definecolor{milcMagenta}{HTML}{E6007E}
\definecolor{milcVino}{HTML}{5A0038}
\definecolor{milcTurquesa}{HTML}{0093A5}
\definecolor{milcVerde}{HTML}{58A923}
\definecolor{milcMostaza}{HTML}{E5B400}
\definecolor{milcOcre}{HTML}{B86600}
\definecolor{milcNegro}{HTML}{241816}
\definecolor{milcGris}{HTML}{F7F7F4}
\definecolor{milcLinea}{HTML}{E8E2DD}
\definecolor{milcGradePrimary}{HTML}{FF6600}
\definecolor{milcGradeDark}{HTML}{9A3E00}
\definecolor{milcGradeSoft}{HTML}{FFE4D0}
\definecolor{milcGradeAccent}{HTML}{CFE7FF}
\definecolor{milcGradeText}{HTML}{FFFFFF}
\color{milcNegro}

\pagestyle{fancy}
\fancyhf{}
% Cabecera de dos lineas: identidad de la guia arriba y, a la derecha y en
% gris, la estacion en curso (\markright desde \milcestacion). Antes de la
% primera estacion la segunda linea va vacia; el \strut mantiene la altura
% para que la linea de arriba no baile entre paginas.
\lhead{\small Tecnología e Informática · Grado 8\\[-1pt]{\footnotesize\strut}}
\rhead{\small Guía 9 · MILC v3\\[-1pt]{\footnotesize\strut\color{milcNegro!65}\rightmark}}
\cfoot{\small\thepage}
\renewcommand{\headrulewidth}{0pt}

% Sobre mostaza (#E5B400) y verde (#58A923) el texto va en milcNegro; sobre el
% resto de la paleta, en blanco. Se usa dentro de fonttitle/fontupper, donde un
% \color posterior manda sobre coltitle.
\newcommand{\milctintasobre}[1]{%
  \ifboolexpr{test{\ifstrequal{#1}{milcMostaza}} or test{\ifstrequal{#1}{milcVerde}}}%
    {\color{milcNegro}}{\color{white}}}

\newtcolorbox{titlebox}[1]{
  enhanced,colback=#1,colframe=#1,arc=7mm,boxrule=0pt,
  left=7mm,right=7mm,top=3mm,bottom=3mm,
  before skip=8pt,after skip=8pt,
  fontupper=\sffamily\bfseries\Large\color{white}
}

\newtcolorbox{softbox}[3]{
  enhanced,breakable,pad at break=3mm,
  colback=#2,colframe=#1,borderline west={4pt}{0pt}{#1},
  arc=2mm,boxrule=.4pt,left=4mm,right=4mm,top=3mm,bottom=3mm,
  before skip=6pt,after skip=8pt,title={#3},coltitle=white,
  colbacktitle=#1,fonttitle=\sffamily\bfseries\milctintasobre{#1},fontupper=\RaggedRight,
  attach boxed title to top left={xshift=3mm,yshift=-2mm},
  boxed title style={arc=2mm, boxrule=0pt}
}

\newtcolorbox{drawbox}[2]{
  enhanced,breakable,pad at break=4mm,
  colback=white,colframe=#1,arc=4mm,boxrule=1pt,
  left=5mm,right=5mm,top=4mm,bottom=4mm,before skip=8pt,after skip=8pt,
  title={#2},coltitle=white,colbacktitle=#1,fonttitle=\sffamily\bfseries\milctintasobre{#1},
  fontupper=\RaggedRight,
  attach boxed title to top left={xshift=4mm,yshift=-2mm},
  boxed title style={arc=3mm, boxrule=0pt}
}

\newtcolorbox{infoband}[1]{
  enhanced,breakable,pad at break=2mm,colback=#1!8,colframe=#1!50,arc=2mm,boxrule=.5pt,
  left=4mm,right=4mm,top=2mm,bottom=2mm,before skip=4pt,after skip=4pt,
  fontupper=\small\RaggedRight
}

% Puente narrativo entre secciones (contrato editorial Conéctate).
% Da continuidad: "Acabas de X. Ahora vas a Y porque Z."
% Visualmente: banda izquierda en color del grado, texto en itálica pequeña.
\newtcolorbox{puentebox}{
  enhanced,breakable,pad at break=2mm,colback=milcGradeSoft,colframe=milcGradeDark,
  borderline west={3pt}{0pt}{milcGradeDark},
  arc=0mm,boxrule=0pt,left=5mm,right=4mm,top=2.5mm,bottom=2.5mm,
  before skip=4pt,after skip=8pt,
  fontupper=\itshape\small\color{milcNegro}\RaggedRight
}
% La flecha va en modo matematico a proposito: Helvetica Neue NO tiene el glifo
% U+2192, asi que \textrightarrow se compilaba a NADA («Missing character: There
% is no → in font Helvetica Neue») y el puente salia sin flecha. En modo
% matematico la flecha viene de la fuente matematica y si se dibuja.
\newcommand{\puente}[1]{%
  \begin{puentebox}%
    \textcolor{milcGradeDark}{$\rightarrow$}\hspace{2mm}#1%
  \end{puentebox}%
}

\newcommand{\linead}[1]{\noindent\color{milcLinea}\rule{#1}{0.5pt}\color{milcNegro}}
\newcommand{\respuesta}[1]{\par\vspace{2mm}\foreach \n in {1,...,#1}{\noindent\color{milcLinea}\rule{\linewidth}{0.5pt}\par\vspace{5mm}}\color{milcNegro}}
\newcommand{\checkbox}{\textcolor{milcGradeDark}{\fbox{\phantom{x}}}\hspace{5pt}}

% ─── Listas aireadas para las guías socioemocionales ────────────────────
% Componer las verificaciones como «\checkbox texto\\» deja el texto que salta
% de línea debajo del cuadrito y apelmaza el bloque. Estos entornos alinean la
% continuación y airean los ítems. Los usa Territorio Interior; cualquier
% familia puede hacerlo.
\newlist{checklista}{itemize}{1}
\setlist[checklista]{
  label=\textcolor{milcGradeDark}{\fbox{\phantom{x}}},
  leftmargin=9mm, labelsep=3.2mm, itemsep=2.4mm,
  topsep=2.5mm, parsep=0pt, partopsep=0pt, before=\RaggedRight
}
\newlist{pasoslista}{enumerate}{1}
\setlist[pasoslista]{
  label=\textcolor{milcGradeDark}{\sffamily\bfseries\arabic*.},
  leftmargin=8mm, labelsep=3mm, itemsep=2.8mm,
  topsep=1mm, parsep=0pt, partopsep=0pt, before=\RaggedRight
}

\newcommand{\phasecard}[4]{
\begin{tcolorbox}[enhanced,colback=#3,colframe=#2,arc=3mm,boxrule=.5pt,height=3.05cm,valign=center,halign=center,left=1.5mm,right=1.5mm,top=2mm,bottom=2mm]
{\sffamily\bfseries\fontsize{22}{24}\selectfont\textcolor{#2}{#1}}\par
{\sffamily\bfseries\small #4}
\end{tcolorbox}
}

% ── Piezas del rediseno 2026 (legibilidad 6.º-11.º) ───────────────────
% Banda de estacion: reemplaza el titlebox macizo. Titulo a la izquierda,
% dato practico (tiempo/modalidad) a la derecha, en una sola linea.
% Nunca huerfana: pide 9 lineas libres (o salta de pagina rellenando con
% \vfill) y prohibe el salto justo despues de la banda. Nueve y no seis:
% la caja partible que sigue necesita ~80pt (titulo adosado, relleno y dos
% lineas) para arrancar en la misma pagina; con menos, tcolorbox la manda
% entera a la siguiente y la banda se queda sola. El penalty 10000 va
% en `after`, ANTES del espacio posterior: un `after skip` (aunque sea 0pt)
% es un \vskip tras la caja, y ese glue es un punto de corte legal que un
% \nopagebreak escrito despues de \end{tcolorbox} ya no cierra. Cada banda
% deja marcador PDF y actualiza la cabecera derecha.
\newcounter{milcestacion}
\newcommand{\milcestacion}[2]{%
\Needspace*{9\baselineskip}%
\stepcounter{milcestacion}%
\pdfbookmark[1]{#2}{milcestacion\themilcestacion}%
\markright{#2}%
\begin{tcolorbox}[enhanced,colback=#1,colframe=#1,arc=2mm,boxrule=0pt,
  left=5mm,right=5mm,top=2.6mm,bottom=2.6mm,before skip=11pt,
  after={\par\nopagebreak\vskip7pt}]
{\sffamily\bfseries\fontsize{15}{18}\selectfont\milctintasobre{#1}#2}
\end{tcolorbox}}

% Tarjeta de paso numerada: sustituye las tablas «Pilar / Aplicacion», que
% eran formato de consulta docente, no de estudiante. El numero va en un
% circulo sobre el borde para que la secuencia se lea de un vistazo.
\newcommand{\milcpaso}[3]{%
\Needspace{4\baselineskip}%
\begin{tcolorbox}[enhanced,colback=white,colframe=milcLinea,arc=1.5mm,boxrule=.9pt,
  left=14mm,right=4mm,top=3mm,bottom=3mm,before skip=4pt,after skip=4pt,
  fontupper=\RaggedRight,
  overlay={\node[anchor=north west,circle,fill=#2,text=white,inner sep=0pt,
    minimum size=7.4mm,font=\sffamily\bfseries\small]
    at ([xshift=3.6mm,yshift=-3.2mm]frame.north west) {#1};}]
#3
\end{tcolorbox}}

% Casilla marcable de tamano util con lapiz (la anterior era \fbox{\phantom{x}},
% de alto variable segun el texto que la seguia).
\newcommand{\milcmarca}{\framebox[3.8mm]{\rule{0pt}{3.4mm}}}

% Fila marcable de lista suelta (checks de la estacion 1).
\newcommand{\milccheckrow}[1]{%
\par\vspace{1.8mm}\noindent
\makebox[6.4mm][l]{\raisebox{-.55mm}{\textcolor{milcNegro}{\milcmarca}}}%
\parbox[t]{\dimexpr\linewidth-6.4mm\relax}{\RaggedRight #1}\par}

% Celda marcable dentro de la rubrica (tabla de autoevaluacion). Misma
% sangria francesa, pero pensada para vivir dentro de una columna X.
\newcommand{\milcopcion}[1]{%
\makebox[5.2mm][l]{\raisebox{-.55mm}{\textcolor{milcNegro}{\milcmarca}}}%
\parbox[t]{\dimexpr\linewidth-5.2mm\relax}{\RaggedRight #1}}

% ── Recursos visuales de la guía ──────────────────────────────────────
% Declarados en el bloque `recursos:` del YAML y emitidos por el builder.
% \guiaFigura   → imagen rasterizada o PDF (foto, carta celeste, montaje).
% \guiaDiagrama → fuente TikZ/circuitikz (.tex) incrustada como vector:
%                 es la vía para circuitos y esquemáticos editables.
% [#1] = texto alternativo (alt) · #2 = ruta relativa al .tex · #3 = pie
% El alt llega al PDF etiquetado (/Alt) y lo lee el lector de pantalla.
\newcommand{\guiaFigura}[3][]{%
\begin{center}
\includegraphics[width=0.88\linewidth,height=0.38\textheight,keepaspectratio,alt={#1}]{#2}\\[1.5mm]
{\footnotesize\itshape\textcolor{milcNegro!75}{#3}}
\end{center}
}

\newcommand{\guiaDiagrama}[2]{%
\begin{center}
\input{#1}\\[1.5mm]
{\footnotesize\itshape\textcolor{milcNegro!75}{#2}}
\end{center}
}

\begin{document}
\thispagestyle{empty}

% ── Portada ───────────────────────────────────────────────────────────
% Antes: un rectangulo de color a pagina completa. Se veia bien en pantalla,
% pero en la fotocopiadora del colegio se comia el toner y salia gris sucio.
% Ahora la banda ocupa el tercio superior y el resto va en blanco.
\begin{tikzpicture}[remember picture,overlay]
  \path[fill=milcGradePrimary]
    (current page.north west) rectangle ([yshift=-10.6cm]current page.north east);

  \node[anchor=north west,text=milcGradeText] at ([xshift=2.0cm,yshift=-1.55cm]current page.north west)
    {\sffamily\bfseries\fontsize{12}{14}\selectfont GRADO};
  \node[anchor=north west,text=milcGradeText,opacity=.85] at ([xshift=2.0cm,yshift=-2.35cm]current page.north west)
    {\sffamily\bfseries\fontsize{8.5}{10}\selectfont SERIE GUÍAS · TECNOLOGÍA E INFORMÁTICA};
  \node[anchor=north east,text=milcGradeText,opacity=.85,align=right] at ([xshift=-2.0cm,yshift=-1.55cm]current page.north east)
    {\sffamily\bfseries\fontsize{9}{12}\selectfont I.E. Sor María Juliana\\Cartago, Valle del Cauca};

  \node[anchor=west,text=milcGradeText] (gradonum) at ([xshift=1.85cm,yshift=-7.1cm]current page.north west)
    {\sffamily\bfseries\fontsize{112}{112}\selectfont 8};
  % El circulito de grado (°) solo aplica a las guias de grado («11°»); en
  % Bebras y Semillero el numero grande es un momento/modulo, no un grado.
  % Se posiciona relativo al nodo `gradonum`, no en coordenadas absolutas.
  \draw[milcGradeText,line width=1.6pt]
    ([xshift=2.0mm,yshift=-4.5mm]gradonum.north east) circle (.15cm);

  \node[anchor=west,text=milcGradeText,text width=11.4cm,align=left] at ([xshift=1.5cm,yshift=0.5cm]gradonum.east)
    {
     {\sffamily\bfseries\fontsize{9}{11}\selectfont\MakeUppercase{Período 1 · Análisis de datos con phronesis}}\\[3mm]
     {\sffamily\bfseries\fontsize{21}{25}\selectfont\setlength{\baselineskip}{27pt}Mini estudio\\[4pt]de datos}\\[3.5mm]
     {\sffamily\bfseries\fontsize{15}{18}\selectfont Guía 9-8-TIC}
    };
\end{tikzpicture}

\vspace*{9.4cm}
\vfill

{\sffamily\bfseries\fontsize{8}{10}\selectfont\textcolor{milcGradeDark}{LO QUE TE LLEVAS HOY}}

\begin{tcolorbox}[enhanced,colback=white,colframe=milcNegro,arc=2mm,boxrule=1.6pt,
  left=5mm,right=5mm,top=4mm,bottom=4mm,before skip=3pt,after skip=12pt,
  fontupper=\RaggedRight\fontsize{11}{15.5}\selectfont]
Un mini estudio en ocho etapas sobre una pregunta real del colegio: pregunta, hipótesis, recolección, limpieza con bitácora, cálculo con las cuatro funciones, gráfico honesto, hallazgo en una frase y decisión. Un informe de dos páginas que documente las ocho etapas, con un párrafo final de limitaciones y de cómo mejorarías el estudio.
\end{tcolorbox}

\begin{tcolorbox}[enhanced,colback=milcGris,colframe=milcGris,arc=2mm,boxrule=0pt,
  left=5mm,right=5mm,top=3.5mm,bottom=3.5mm,after skip=12pt]
{\sffamily\bfseries\fontsize{8}{10}\selectfont\textcolor{milcNegro!55}{ESTA GUÍA ES DE}}\\[3mm]
\begin{tabularx}{\linewidth}{X p{3.2cm} p{3.2cm}}
\linead{\linewidth} & \linead{3.0cm} & \linead{3.0cm}\\[-1mm]
{\fontsize{8.5}{10}\selectfont\textcolor{milcNegro!55}{Nombre completo}} &
{\fontsize{8.5}{10}\selectfont\textcolor{milcNegro!55}{Grupo}} &
{\fontsize{8.5}{10}\selectfont\textcolor{milcNegro!55}{Fecha}}\\
\end{tabularx}
\end{tcolorbox}

\vfill

\noindent\textcolor{milcLinea}{\rule{\linewidth}{1pt}}\\[2mm]
{\sffamily\bfseries\fontsize{9.5}{12}\selectfont Autor: Dr. Álvaro Cárdenas Orozco}\\[1mm]
{\sffamily\fontsize{8.5}{11}\selectfont\textcolor{milcNegro!55}{Metodología MILC · apertura ancestral · pensamiento computacional · triángulo Dussel–estoicismo–Floridi}}

\newpage

\section*{Apertura: Mini estudio de datos --- ocho etapas para responder una pregunta del colegio}

\begin{softbox}{milcGradeDark}{milcGradeSoft}{Saber ancestral}
Cuando una comunidad indígena pide constituir o ampliar su resguardo, se levanta un \textbf{censo}. El instructivo oficial lo define como «una fotografía de la realidad de una comunidad en un tiempo y espacio específico» (Agencia Nacional de Tierras, 2023). No lo hace un funcionario solo. Se apoya en miembros del cabildo, maestros y líderes de la comunidad, que van casa por casa con un formato. Se pregunta a todas las personas: quiénes son, cuántos años tienen, qué hacen. Después alguien verifica lo escrito, lo sistematiza y lo cuenta. Y con esa lista se toma una decisión enorme: la tierra se distribuye «conforme al censo realizado». Mira las etapas: una pregunta, un método, una recolección persona por persona, una revisión, un cálculo, una decisión. Es un estudio completo. La \textbf{cara de exclusión} es literal: quien no queda en la lista el día del censo no existe para el reparto. Hoy vas a hacer un estudio mucho más pequeño, sobre tu colegio, con las mismas etapas y el mismo cuidado por quién queda fuera.\par\smallskip{\footnotesize\textcolor{milcNegro!70}{\textbf{Fuente.} Agencia Nacional de Tierras, Dirección de Asuntos Étnicos. (2023). \emph{Instructivo: levantamiento y diligenciamiento del formato de censo poblacional para comunidades indígenas} (ACCTI-I-021, versión 1).}}
\end{softbox}

\begin{softbox}{milcTurquesa}{milcGris}{Saber contemporáneo}
Un \textbf{mini estudio de datos} usa todo lo que aprendiste en el periodo para responder una pregunta real. Tiene ocho etapas. (1) \textbf{Pregunta}: una sola, clara, que se pueda responder con datos. (2) \textbf{Hipótesis}: lo que crees que vas a encontrar, escrito antes de mirar. (3) \textbf{Recolección}: cómo consigues los datos, cuántos y cuándo. (4) \textbf{Limpieza} con bitácora, como en la sesión 3. (5) \textbf{Cálculo}: las cuatro funciones y una fórmula compuesta. (6) \textbf{Gráfico honesto}, como en la sesión 7. (7) \textbf{Hallazgo}: lo más importante, en una frase. (8) \textbf{Decisión}: qué se haría con eso. Saltar una etapa rompe el estudio. Y la última página siempre dice qué no pudiste saber.
\end{softbox}

\begin{softbox}{milcMostaza}{milcGris}{Pregunta-puente}
\textit{El censo del resguardo decide la tierra, y quien no está en la lista queda fuera. Cuando encuestes a veinte compañeros sobre su tiempo de pantalla, ¿quién no va a estar en tu lista? ¿Y qué le pasa a tu hallazgo por esa ausencia?}
\end{softbox}

\begin{softbox}{milcVerde}{milcGris}{Saber y saber hacer}
Al terminar podrás: (1) \textbf{identificar} preguntas del colegio que se puedan responder con datos en una semana; (2) \textbf{evaluar} cuál de ellas vale la pena con cuatro filtros; (3) \textbf{crear} un mini estudio completo de ocho etapas; (4) \textbf{explicar} sus limitaciones y cómo lo mejorarías.
\end{softbox}



\small
\begin{tabularx}{\textwidth}{>{\bfseries}p{3.0cm}Y}
\rowcolor{milcGradePrimary!12}Autor & Dr. Álvaro Cárdenas Orozco\\
DBA & Tratamiento de datos cuantitativos con criterio ético (MEN, T\&I)\\
Referentes & Phronesis aristotélica · Floridi (ética del dato) · Estoicismo (ver lo que es)\\
Producto final & Un mini estudio en ocho etapas sobre una pregunta real del colegio: pregunta, hipótesis, recolección, limpieza con bitácora, cálculo con las cuatro funciones, gráfico honesto, hallazgo en una frase y decisión. Un informe de dos páginas que documente las ocho etapas, con un párrafo final de limitaciones y de cómo mejorarías el estudio.\\
\end{tabularx}
\normalsize

\puente{En la sesión 8 tu hoja aprendió a cuidarse sola. Hoy juntas todo el periodo en un estudio propio. En la sesión 10 lo vas a sustentar en cuatro minutos frente al grupo.}

\milcestacion{milcGradeDark}{Ruta MILC: así vamos a aprender}
\begin{tabularx}{\textwidth}{>{\centering\arraybackslash}X>{\centering\arraybackslash}X>{\centering\arraybackslash}X>{\centering\arraybackslash}X}
\phasecard{1}{milcVerde}{milcGris}{Escucha\\[-1mm]\small Observo, escucho y pregunto.} &
\phasecard{2}{milcTurquesa}{milcGris}{Entender\\[-1mm]\small Sistematizo y ordeno.} &
\phasecard{3}{milcMagenta}{milcGris}{Hacer\\[-1mm]\small Construyo y aplico.} &
\phasecard{4}{milcVino}{milcGris}{Cierre\\[-1mm]\small Evalúo y reflexiono.}
\end{tabularx}


\puente{Antes de elegir, vas a escribir diez preguntas posibles sobre tu colegio. Diez, no una: la buena pregunta casi nunca es la primera.}

\milcestacion{milcVerde}{Estación 1 de 3 · Escucha}

\begin{softbox}{milcVerde}{milcGris}{Escena para escuchar}
\textbf{Actividad 1 · IDENTIFICA --- Diez preguntas posibles} (15 min · individual). (1) Escribe diez preguntas sobre tu colegio que se puedan responder con datos. Por ejemplo: ¿cuántos minutos de pantalla al día tienen los de octavo? ¿Qué producto se vende más en la tienda cada día? ¿Cuántos vienen a pie, en bus o en bicicleta? (2) Cada pregunta debe tener algo que se cuente o se mida. (3) Tacha las que no podrías responder en una semana. (4) Marca con una estrella las tres que más te interesan de verdad.
\end{softbox}

{\sffamily\bfseries\fontsize{8}{10}\selectfont\textcolor{milcNegro!55}{MARCA CUANDO LO HAYAS HECHO}}
\milccheckrow{Tengo diez preguntas escritas y cada una tiene algo que se cuenta o se mide.}
\milccheckrow{Taché las que no se pueden responder en una semana.}
\milccheckrow{Marqué con estrella las tres que más me interesan.}

\begin{infoband}{milcVerde}


\textbf{Lo que va al cuaderno (Actividad 1 · IDENTIFICA).} \textbf{Título:} «Actividad 1 --- Diez preguntas posibles». \textbf{Formato:} lista numerada de diez preguntas, con tachones y estrellas. \textbf{Extensión:} media página. \textbf{Sabes que terminaste cuando} quedan al menos cinco sin tachar y tres tienen estrella.
\end{infoband}

\puente{Ya tienes diez preguntas. Ahora las pasas por cuatro filtros con tu pareja, eliges una, escribes la hipótesis antes de mirar nada y planeas la recolección.}

\milcestacion{milcTurquesa}{Estación 2 de 3 · Entender}

\begin{softbox}{milcTurquesa}{milcGris}{Lo que vas a entender}
\textbf{Actividad 2 · EVALÚA --- Elige la pregunta y escribe la hipótesis} (25 min · parejas). (1) Con tu pareja, pasen sus tres preguntas con estrella por cuatro filtros: ¿consigo los datos esta semana?, ¿voy a tener al menos quince datos?, ¿de verdad quiero saber la respuesta?, ¿el hallazgo serviría para algo? (2) Elijan cada uno la pregunta que pase los cuatro. (3) Escriban la hipótesis antes de recolectar nada: «creo que voy a encontrar X, porque Y». (4) Planeen la recolección: a quién, cuántos, cuándo y con qué, encuesta, conteo u observación. (5) Escriban quién va a quedar fuera de la muestra y por qué.
\end{softbox}

\milcpaso{1}{milcTurquesa}{\textbf{Los cuatro filtros.} Datos accesibles esta semana. Muestra de al menos quince. Interés real. Algún uso para el hallazgo. La pregunta que pasa los cuatro es la tuya; si ninguna pasa, vuelve a la lista de diez.}
\milcpaso{2}{milcTurquesa}{\textbf{La hipótesis va antes.} Si la escribes después de mirar los datos, cualquier resultado parece esperado. Escribirla antes te permite equivocarte, y equivocarte con datos es aprender algo. Decir «mi hipótesis estaba mal» es una de las frases más honestas de un estudio.}
\milcpaso{3}{milcTurquesa}{\textbf{Quién queda fuera.} Si encuestas a tus amigos, tu muestra son tus amigos. Si solo observas los lunes, tu estudio es de lunes. Esa no es una falla: es una limitación, y se declara en la última página.}
\milcpaso{4}{milcTurquesa}{\textbf{Las ocho etapas, en orden.} Pregunta. Hipótesis. Recolección. Limpieza con bitácora. Cálculo con las cuatro funciones y una fórmula. Gráfico honesto. Hallazgo en una frase. Decisión. Ninguna se salta.}

\begin{drawbox}{milcTurquesa}{Las ocho etapas, en una mirada}
\textbf{1. Pregunta:} una sola, medible. \\[1mm] \textbf{2. Hipótesis:} antes de mirar. \\[1mm] \textbf{3. Recolección:} a quién, cuántos, cuándo. \\[1mm] \textbf{4. Limpieza:} con bitácora. \\[1mm] \textbf{5. Cálculo:} las cuatro funciones. \\[1mm] \textbf{6. Gráfico:} honesto. \\[1mm] \textbf{7. Hallazgo:} una frase. \\[1mm] \textbf{8. Decisión:} y limitaciones.
\end{drawbox}

\begin{softbox}{milcMostaza}{milcGris}{Errores comunes y su corrección}
(1) \textbf{Pregunta vaga}: «¿cómo está el colegio?» no se responde con datos. (2) \textbf{Hipótesis después}: todo parece esperado. (3) \textbf{Muestra de cinco}: no alcanza para concluir nada. (4) \textbf{Cálculos de adorno}: fórmulas que no responden la pregunta. (5) \textbf{Hallazgo inflado}: concluir más de lo que veinte datos soportan. (6) \textbf{Sin limitaciones}: presentar el estudio como definitivo.
\end{softbox}

\begin{infoband}{milcTurquesa}


\textbf{Lo que va al cuaderno (Actividad 2 · EVALÚA).} \textbf{Título:} «Actividad 2 --- Mi pregunta y mi hipótesis». \textbf{Formato:} la pregunta elegida, los cuatro filtros con su respuesta, la hipótesis con su porqué, el plan de recolección en cuatro líneas y una línea sobre quién queda fuera. \textbf{Extensión:} media página. \textbf{Sabes que terminaste cuando} la pregunta pasa los cuatro filtros, la hipótesis está escrita antes de recolectar y sabes a quién no vas a encuestar.
\end{infoband}

\puente{Ya tienes pregunta, hipótesis y plan. Toca ejecutar las ocho etapas durante la semana y armar el informe de dos páginas.}

\milcestacion{milcMagenta}{Estación 3 de 3 · Hacer}

\begin{softbox}{milcMagenta}{milcGris}{Cómo vas a construir tu producto}
\textbf{Actividad 3 · CREA --- El mini estudio completo} (30 min · individual). Hoy dejas listas las primeras etapas; el resto se hace durante la semana. (1) Recolecta los datos durante tres a cinco días, como planeaste. (2) Límpialos con bitácora, como en la sesión 3. (3) Calcula las cuatro funciones y una fórmula compuesta que responda tu pregunta. (4) Haz un gráfico honesto del tipo correcto. (5) Escribe el hallazgo en una frase y la decisión que tomarías. (6) Arma el informe de dos páginas con las ocho etapas y cierra con limitaciones y mejoras.
\end{softbox}

\milcpaso{1}{milcMagenta}{\textbf{El informe tiene ocho secciones.} Una por etapa, de uno a tres párrafos cada una. La pregunta y por qué importa. La hipótesis con su porqué. El método con cuántos datos, cómo y cuándo. La limpieza con el resumen de la bitácora. Las cifras clave. El gráfico con dos líneas de lectura. El hallazgo en una frase. La decisión.}
\milcpaso{2}{milcMagenta}{\textbf{La última página dice lo que no sabes.} Tres a cinco limitaciones concretas: muestra pequeña, solo amigos, solo lunes, datos de memoria. Tres mejoras si lo hicieras de nuevo. Y qué seguirías investigando con más tiempo.}
\milcpaso{3}{milcMagenta}{\textbf{El hallazgo cabe en una frase.} «Los de octavo pasan en promedio 3 horas al día en pantalla, y el máximo fue de 8.» Si necesitas un párrafo para decirlo, todavía no lo encontraste.}
\milcpaso{4}{milcMagenta}{\textbf{Si la hipótesis falló, mejor.} Un estudio que confirma lo que ya creías enseña poco. Uno que te contradice te obliga a mirar de nuevo, y eso es lo que vas a sustentar en la sesión 10.}

\begin{drawbox}{milcMagenta}{Antes de entregar}
\checkbox Pregunta clara y medible, una sola. \\[1mm] \checkbox Hipótesis escrita antes de recolectar. \\[1mm] \checkbox Al menos quince datos, con método explícito. \\[1mm] \checkbox Limpieza con bitácora. \\[1mm] \checkbox Cuatro funciones y una fórmula compuesta. \\[1mm] \checkbox Un gráfico honesto del tipo correcto. \\[1mm] \checkbox Hallazgo en una frase y decisión. \\[1mm] \checkbox Limitaciones y mejoras en la última página.
\end{drawbox}

\begin{softbox}{milcMagenta}{milcGris}{Plantilla de trabajo}
\textbf{Pregunta.} \textit{Una sola, medible.} \\[1mm] \textbf{Hipótesis.} \textit{Creo que…, porque…} \\[1mm] \textbf{Método.} \textit{A quién, cuántos, cuándo, con qué.} \\[1mm] \textbf{Cifras y gráfico.} \textit{Las cuatro funciones, una fórmula, un gráfico.} \\[1mm] \textbf{Hallazgo y decisión.} \textit{Una frase y una acción.} \\[1mm] \textbf{Limitaciones.} \textit{Quién quedó fuera y cómo lo mejoraría.}
\end{softbox}

\begin{infoband}{milcMagenta}


\textbf{Lo que va al cuaderno (Actividad 3 · CREA).} \textbf{Título:} «Actividad 3 --- Mi mini estudio». \textbf{Formato:} el esquema del informe con las ocho secciones tituladas y, debajo de cada una, qué vas a poner; el informe completo se entrega aparte. \textbf{Extensión:} una página. \textbf{Sabes que terminaste cuando} las etapas 1 a 3 están listas, las ocho secciones tienen título y sabes qué te falta recolectar esta semana.
\end{infoband}

\puente{Lo que entregas: el informe de dos páginas con las ocho etapas y el párrafo de limitaciones. La prueba de calidad: tu docente lee la pregunta y el hallazgo, y el hallazgo responde exactamente esa pregunta.}

\milcestacion{milcMostaza}{Producto de la sesión}
\textbf{Mini estudio de ocho etapas + informe de dos páginas con limitaciones}.
\begin{softbox}{milcMostaza}{milcGris}{Criterios de calidad}
(1) La pregunta es una sola, clara y medible, y pasa los cuatro filtros. (2) La hipótesis está escrita antes de recolectar y dice por qué. (3) Hay al menos quince datos, con método explícito y limpieza con bitácora. (4) El cálculo usa las cuatro funciones y una fórmula compuesta que responden la pregunta. (5) El gráfico es honesto y el hallazgo cabe en una frase. (6) La última página declara limitaciones, mejoras y quién quedó fuera.
\end{softbox}

\puente{Antes de cerrar, mira lo que hiciste desde cinco lados. Un estudio honesto dice a quién no encuestó. Uno deshonesto lo esconde.}

\milcestacion{milcVino}{Evaluación liberadora · cinco dimensiones}

\begin{softbox}{milcVino}{milcGris}{Sentido de la evaluación}
Evaluar no es solo revisar una nota. Es reconocer qué aprendí, qué cambió en mí, qué emociones manejé, cómo me relaciono con la ciudad y la comunidad, y qué le dejaré a quien viene detrás.
\end{softbox}

\textbf{Mi reflexión liberadora en cinco dimensiones.} Completa cada frase iniciada:

\small
\begin{tabularx}{\textwidth}{>{\bfseries\RaggedRight\arraybackslash}p{3.4cm}Y>{\RaggedRight\arraybackslash}p{4.6cm}}
\rowcolor{milcVino}\color{white}Dimensión & \color{white}\bfseries Mi respuesta (frame starter) & \color{white}\bfseries Algo que descubrí\\
1. Personal & Lo que más cambió en mí fue \linead{6.5cm} & \linead{4.2cm}\\
2. Emocional & La emoción más fuerte fue \linead{2.5cm} y la regulé \linead{3cm} & \linead{4.2cm}\\
3. Ciudadana & En democracia esto importa porque \linead{6cm} & \linead{4.2cm}\\
4. Local (Cartago) & En mi barrio sería útil para \linead{6.5cm} & \linead{4.2cm}\\
5. Intergeneracional & A mi abuela/abuelo le diría \linead{6.5cm} & \linead{4.2cm}\\
\end{tabularx}
\normalsize

\begin{infoband}{milcVino}
\textbf{Para la próxima sustentación.} Vuelve a esta tabla en seis meses. Verás que las respuestas cambian — no porque lo aprendido cambie, sino porque tú cambias. La evaluación liberadora es una conversación contigo a través del tiempo.
\end{infoband}


\puente{Para terminar, tres ideas sobre investigar de verdad, contadas con palabras de esta guía: Dussel sobre saber perder tiempo para elegir lo que importa, Marco Aurelio sobre cambiar de idea cuando te muestran el error, Floridi sobre saber qué preguntar antes de buscar datos.}

\milcestacion{milcGradeDark}{Tres ideas para cerrar · Dussel, un estoico y Floridi}

\begin{softbox}{milcVino}{milcGris}{Enrique Dussel · lente del nosotros}
\textit{Como los temas son infinitos y el tiempo corto, hay que saber perder tiempo para elegir los temas fundamentales.}

{\footnotesize\itshape\color{milcNegro!70}--- idea tomada de Enrique Dussel, contada con palabras de esta guía. No es una cita textual.}

Escribir diez preguntas para quedarte con una parece perder tiempo. Es lo contrario: es el tiempo que te evita pasar una semana recolectando datos para una pregunta que no importaba.

\textbf{¿Elegí mi pregunta porque importa, o porque era la más fácil de responder?}
\end{softbox}
\linead{\linewidth}\\[1mm]
\linead{\linewidth}

\begin{softbox}{milcMostaza}{milcGris}{Estoicismo · Marco Aurelio · cuidado interior}
\textit{Si alguien te muestra que estás equivocado, cambia de camino con gusto. Lo que daña no es el error: es aferrarse a él.}

{\footnotesize\itshape\color{milcNegro!70}--- idea tomada de Marco Aurelio, contada con palabras de esta guía. No es una cita textual.}

Tu hipótesis puede fallar. Los datos son «alguien» que te muestra el error. Cambiar de idea con datos no es perder: es exactamente para lo que hiciste el estudio.

\textbf{Si mis datos contradicen mi hipótesis, ¿voy a decirlo en el informe?}
\end{softbox}
\linead{\linewidth}\\[1mm]
\linead{\linewidth}

\begin{softbox}{milcTurquesa}{milcGris}{Luciano Floridi · lente de la infoesfera}
\textit{Ganan quienes saben preguntar y responder, y por eso saben qué datos buscar.}

{\footnotesize\itshape\color{milcNegro!70}--- idea tomada de Luciano Floridi, contada con palabras de esta guía. No es una cita textual.}

Sin pregunta clara, recolectas de todo y no respondes nada. La pregunta decide qué datos, a quién, cuántos. Por eso la etapa uno se hace despacio y las otras siete salen de ella.

\textbf{¿Mi pregunta me dijo exactamente qué datos buscar, o busqué y después pregunté?}
\end{softbox}
\linead{\linewidth}\\[1mm]
\linead{\linewidth}

\begin{infoband}{milcNegro}
\textbf{Nota para el docente.} Las tres ideas están contadas con palabras de esta guía a partir del pensamiento de cada autor; no son citas textuales. De 9.º en adelante se trabajan con la obra y su referencia.
\end{infoband}

\milcestacion{milcVino}{Mi autoevaluación · marca una casilla por fila}

{\small
\setlength{\extrarowheight}{2.2pt}
\begin{tabularx}{\textwidth}{>{\RaggedRight\arraybackslash\bfseries}p{3.0cm}YYY}
\rowcolor{milcVino}\color{white}\bfseries Criterio & \color{white}\bfseries Lo logré & \color{white}\bfseries En proceso & \color{white}\bfseries Necesito apoyo\\[.6mm]
Comprensión del tema & \milcopcion{Explico las ideas con mis palabras.} & \milcopcion{Reconozco las ideas, falta precisión.} & \milcopcion{Necesito repasar lo básico.}\\[.8mm]
\rowcolor{milcGris}Pensamiento computacional & \milcopcion{Usé los pilares al armar mi producto.} & \milcopcion{Usé algunos pilares.} & \milcopcion{Necesito releer la estación 2.}\\[.8mm]
Aplicación y producto & \milcopcion{Mi producto cumple los criterios.} & \milcopcion{Está armado, con ajustes pendientes.} & \milcopcion{Aún está en construcción.}\\[.8mm]
\rowcolor{milcGris}Reflexión 5 dimensiones & \milcopcion{Respondí cada una con honestidad.} & \milcopcion{Respondí algunas con detalle.} & \milcopcion{Mis respuestas son muy generales.}\\[.8mm]
Triángulo de pensamiento & \milcopcion{Conecté las 3 voces con mi trabajo.} & \milcopcion{Conecté al menos una.} & \milcopcion{No las relaciono con mi proceso.}\\[.8mm]
\end{tabularx}}

\vspace{3mm}
\textbf{Hoy entendí que:}\\[1mm]
\linead{\linewidth}\\[1mm]
\linead{\linewidth}

\textbf{Mi compromiso para la próxima sesión es:}\\[1mm]
\linead{\linewidth}\\[1mm]
\linead{\linewidth}

\begin{softbox}{milcMostaza}{milcGris}{Cita de despedida · Marco Aurelio}
\textit{``El obstáculo en el camino se vuelve el camino. Lo que estorba al camino se vuelve camino.''}\\[1mm]
Cada error de hoy, bien observado, se convierte en pieza del camino que pisarás mañana.
\end{softbox}

\small
\begin{tabularx}{\textwidth}{>{\bfseries}p{3.6cm}Y}
\rowcolor{milcGradeDark!12}Nombre completo: & \linead{10cm}\\
Fecha: & \linead{4cm} \quad Curso/grupo: \linead{4cm}\\
Mi firma: & \linead{9cm}\\
\end{tabularx}
\normalsize

\begin{infoband}{milcGradeDark}
\textbf{Para la próxima sesión.} Llega con el mini estudio terminado y el informe de dos páginas. En la sesión 10 lo vas a sustentar en cuatro minutos frente al grupo, con cinco diapositivas, y vas a responder preguntas. Un estudio honesto se sustenta solo; uno inflado se cae en la primera pregunta.
\end{infoband}

\end{document}
